%%% Introduction / Motivation

Whole-genome alignments are a foundational data layer for pangenomics \cite{computational2018, Eizenga2020b, Paten2017}, yet current workflows treat them as intermediates consumed during graph construction \cite{Garrison2024, Hickey2023}. We propose that the alignment collection itself---together with the transitive relationships it encodes---constitutes a directly queryable pangenome representation. Operating in this alignment space, we can project coordinates transitively across genomes, partition the pangenome along alignment connectivity, and materialize explicit graphs only when needed for specific regions of interest. As high-quality assemblies are now routinely produced for hundreds to thousands of individuals \cite{Liao2023, Wang2022}, this alignment-first approach avoids the substantial computation required for monolithic graph construction at population scale.

Pairwise alignments inherently define a variation graph: positions matched transitively across alignments correspond to shared nodes in a graph that losslessly encodes both the sequences and their relationships \cite{garrison2023seqwish}. The pangenome graph is thus implied by the alignments, even when it is never constructed. Graph construction methods---seqwish \cite{garrison2023seqwish}, Cactus \cite{Paten2011cactus}, TBA/MULTIZ \cite{Blanchette2004}---exploit this property to build pangenome graphs and multiple alignments, but require materializing the complete structure before any region can be examined. Coordinate projection tools such as liftOver \cite{Kuhn2013ucsc} and paftools \cite{Li2018minimap2} operate at query time but map positions between only two genomes; halLiftover \cite{Hickey2013hal} achieves multi-hop projection but requires a pre-computed hierarchical alignment. No existing tool exploits this transitive structure at query time directly from pairwise alignments, without first constructing a graph or hierarchical alignment.
